\chapter{問いと答え}\label{ch:qa}

この章では、研究の問いと、2 週間で得た答えを、根拠の数値とともに 1 か所にまとめる。各結果の詳しい説明は、括弧の中の章にある。

\section{出発点：元論文}

出発点は、Lee・Kim・Kim・Lee の論文（arXiv:2504.14345、ICLR 2025 の Financial AI のワークショップ）とその公式実装である。手法は次の 3 段階でできている。

\begin{enumerate}
\item 各銘柄について、直近 1 か月の日次リターンの列と会社の情報を LLM に渡し、\textbf{翌月の平均日次リターン}を 1 つの数で答えさせる。
\item 同じ質問をくり返し、回答の平均をビュー $Q$、回答の分散をビューの不確かさ $\Omega$ にする。
\item 市場の均衡リターンとビューを Black--Litterman モデルで混ぜ、平均分散最適化で投資比率を決める。
\end{enumerate}

論文は、この方法で作ったポートフォリオが等金額や平均分散最適化を上回ったと報告している。私たちはこれを OpenAI の gpt-4o-mini で、\SP{} の全構成銘柄について 30 か月分再現し、何が起きているのかを調べるところから始めた。

\section{問いの立て方の変化}

研究の問いは、結果を受けて 3 回変わった（\figref{fig:qa-questions}）。

\begin{figure}[htbp]
\centering
\begin{tikzpicture}[node distance=5mm and 6mm]
\node[box=c1,text width=32mm] (q1) {\textbf{問い 1}\\LLM の数値予想は\\BLM のビューになるか};
\node[box=c2,text width=32mm,right=of q1] (q2) {\textbf{問い 2}\\LLM は文章から\\株価の材料を読めるか};
\node[box=c3,text width=32mm,right=of q2] (q3) {\textbf{問い 3}\\読めた情報で\\運用成績は上がるか};
\node[box=c7,text width=32mm,right=of q3] (q4) {\textbf{問い 4}\\どこを直せば\\良くなるか};
\draw[arr] (q1) -- (q2); \draw[arr] (q2) -- (q3); \draw[arr] (q3) -- (q4);
\node[lab,below=2mm of q1,text width=34mm,align=center] {総合報告書 1、実験 1\\（10/9〜17）};
\node[lab,below=2mm of q2,text width=34mm,align=center] {実験 2〜5・10\\（10/17〜19）};
\node[lab,below=2mm of q3,text width=34mm,align=center] {実験 7〜9、X01〜X05\\（10/17〜19）};
\node[lab,below=2mm of q4,text width=34mm,align=center] {実験 11\\（10/19〜21）};
\end{tikzpicture}
\caption{問いの移り変わり。}
\label{fig:qa-questions}
\end{figure}

\begin{description}[leftmargin=1.5em,style=nextline,itemsep=0.5em]
\item[問い 1：LLM の数値予想は BLM のビューになるか] 元論文の形のまま調べた。答えは「ならない」。LLM の予想は入力の値動きをなぞるだけで、それ以外の予測力は検出できなかった。
\item[問い 2：LLM は文章から株価の材料を読めるか] 数字ではなく、決算発表文という\textbf{文章}を読ませ、その点数が\textbf{発表への株価の反応}を説明するかを調べた。答えは「読める」。数値の驚き、アナリスト予想との差、従来の文章分析、反応で学習したモデルを統制しても残った。
\item[問い 3：読めた情報で運用成績は上がるか] 読めた情報を月次や日次の BLM に入れた。答えは「上がらない」。情報は発表日のうちに株価に入りきっていた。
\item[問い 4：どこを直せば良くなるか] 文章の読ませ方（掃除、読む範囲、埋め込みの取り方、圧縮）と、点数のまとめ方を機械的な規則で試した。答えは「ほとんど直らない」。ただし先読みのない LLM が従来の NLP を超えることは、3 つの期間の切り方すべてで確かめられた。
\end{description}

\section{答えの一覧}

\begin{longtable}{P{2.3cm}P{6.6cm}P{5.6cm}}
\caption{主な結果。「支持」は、事前に決めた判定の規則で支持されたこと。}\label{tab:qa-main}\\
\toprule
\textbf{実験} & \textbf{問い} & \textbf{結果} \\
\midrule
\endfirsthead
\toprule
\textbf{実験} & \textbf{問い} & \textbf{結果} \\
\midrule
\endhead
\bottomrule
\endfoot
\multicolumn{3}{l}{\textbf{問い 1：数値予想}} \\
実行 1 & 元論文と同じ聞き方で予想を集める & 予想と入力の月の平均リターンの相関 0.89。IC は約 +0.05 \\
最初の BLM & 元論文の計算を再現 & 事後の期待リターンの 99.9\% がビューの側。最適化は途中で止まっていた \\
実験 A〜G & 入力と質問の形を 7 通りに変える & どの形でも、答えは直近の値動きに沿った \\
H02 & LLM に固有の予測力 & IC $-0.004$（95\% 区間の上限 +0.027）。\textbf{なし} \\
H04 & 校正した $\Omega$ の BLM & 置換の対照と区別できない（$p=0.104$） \\
H05 & 文章だけを渡す & なぞりは消えた（相関 +0.77 → +0.05）が、当たるかは不明 \\
H06 & 決算発表文の点数 → 翌月のリターン & +0.017（$t$ 0.97）。月次では支持されない \\
\midrule
\multicolumn{3}{l}{\textbf{問い 2：文章と株価の反応（発表をはさむ 2 取引日）}} \\
実験 2 & gpt-4o-mini の点数 & \textbf{+120bp／標準偏差}（$t$ 9.41）。支持 \\
実験 3 & アナリスト予想との差を統制 & \textbf{+116bp}（$t$ 9.37）。97\% が残る。支持 \\
実験 4 & 反応で追加学習した FinBERT と同時 & \textbf{+102bp}（$t$ 7.77）。説明力は LLM が約 2.4 倍。支持 \\
実験 5 & 先読みのない ChronoGPT の埋め込み＋リッジ（2016〜2023） & \textbf{+33.6bp}（$t$ 5.52）。TF-IDF を吸収。支持 \\
実験 10 & ChronoGPT を LoRA で追加学習 & +29.5bp（$t$ 2.40）。FinBERT と同時で +0.4bp \\
実験 11 & 先読みのない ChronoGPT が NLP 3 つを超えるか（2020〜2026） & \textbf{+34bp}（$t$ 4.79）。3 つの切り方すべてで支持 \\
\midrule
\multicolumn{3}{l}{\textbf{問い 3：運用成績}} \\
実験 2・3 & 発表の後のドリフト（2〜21 日目） & +10bp（$t$ 0.83）。有意でない \\
実験 8 & 発表の翌日（day 1） & +3.6bp（$t$ 0.93）。支持されない \\
実験 7 & 月次の BLM（31 か月） & 費用の後 +0.10\%／月（$t$ 1.22）。支持されない \\
実験 9 & 月次の BLM（約 10 年、先読みなし） & +0.003\%／月（$t$ 0.04）。支持されない \\
実験 8 & 日次の BLM & 費用の後 年 −20\%（回転率 1 日 42\%） \\
\midrule
\multicolumn{3}{l}{\textbf{問い 4：直し方}} \\
実験 11 & 修正 1〜5（掃除、読む範囲、まとめ方、埋め込み、圧縮） & 開発の規則で残った修正はほぼなし \\
実験 11（探索） & 中間の層の埋め込み & R$^2$ 増分 0.49 → 0.80（$t$ 2.1〜2.9）。今後の候補 \\
\midrule
\multicolumn{3}{l}{\textbf{ポートフォリオの道具}} \\
X01 & 共分散のノイズ除去 & BLM の最適化では効かない（Ledoit--Wolf で十分） \\
X02 & 共分散の窓・EWMA & 実現の TE は上限の約 1.4 倍のまま \\
X03 & TE の上限の校正 & 実現の TE 4.79\% → 3.38\%。約束にほぼ近づく \\
X04 & VIX で共分散を合わせる & 改善しない \\
X05 & $\alpha$ の大きさ & TE の上限が効くので、大きさは結果に関係しない \\
\end{longtable}

\section{答えの形：反応は説明するが、儲からない}

結果を 1 枚の図にすると\figref{fig:qa-chain}のようになる。決算発表文から株価までの道のりを、左から順に確かめた。

\begin{figure}[htbp]
\centering
\begin{tikzpicture}[node distance=4mm and 5mm]
\node[io=24mm] (t) {決算発表文\\（EX-99.1）};
\node[md=24mm,right=of t] (m) {LLM\\gpt-4o-mini／\\ChronoGPT};
\node[io=24mm,right=of m] (s) {点数\\（発表ごと）};
\node[box=c3,text width=28mm,right=of s] (r) {発表をはさむ\\2 日の反応\\\textbf{説明する}};
\node[box=c2,text width=28mm,below=9mm of r] (d) {翌日・翌月の\\リターン\\\textbf{説明しない}};
\node[box=c2,text width=28mm,left=of d] (p) {BLM の\\ポートフォリオ\\\textbf{上回らない}};
\draw[arr] (t) -- (m); \draw[arr] (m) -- (s); \draw[arr] (s) -- (r);
\draw[arr] (r) -- node[lab,right,text width=24mm] {株価に\\入りきる} (d);
\draw[arr] (d) -- (p);
\end{tikzpicture}
\caption{この研究の結論の形。LLM の読み取りは本物だが、市場も同じ日に同じことを読む。}
\label{fig:qa-chain}
\end{figure}

この形は、ファイナンスの教科書の「効率的市場」の考え方とよく合う。公開された文章の情報は、公開と同時に価格に入る。LLM が読めることは、市場の参加者（人間のアナリストやアルゴリズム）もすでに読んでいる。LLM の読み取りが速さで勝つ可能性はあるが、この研究のデータ（日次の終値）では、発表の当日より細かい時間は見られない。

\section{この研究で言えること・言えないこと}

\begin{longtable}{P{7.2cm}P{7.2cm}}
\toprule
\textbf{言えること} & \textbf{言えないこと} \\
\midrule
\endhead
\bottomrule
\endfoot
元論文の計算は、事後の期待リターンがほぼビューだけになり、最適化も途中で止まる。結果は計算の作りに左右される & 元論文の主張が誤りだとは言えない（モデルと期間が違う） \\
gpt-4o-mini に数値を予想させると、入力の直近の値動きをなぞる & どの LLM でもなぞるとは言えない（gpt-4o-mini だけを調べた） \\
LLM の決算発表文の点数は、発表への 2 日の反応を強く説明し、数値の驚き・アナリスト予想・従来の NLP・反応で学習した FinBERT を超える & 発表の当日のうちの、より細かい時間の動きはわからない \\
先読みのない LLM（ChronoGPT）でも、2016〜2026 年の反応を説明する（gpt-4o-mini の結果が記憶ではない傍証） & gpt-4o-mini の点数に記憶が混ざっていないことの証明にはならない \\
その情報は、翌日・翌月のリターンには有意に残らない。月次・日次の BLM は費用の後で時価総額を有意に上回らない & 効果がまったくないとは言えない（区間は小さな正の値も含む） \\
読ませ方とまとめ方の単純な修正は、説明力をほとんど上げない & もっと強い LLM や別の聞き方で上がらないとは言えない \\
\end{longtable}

\section{この研究の貢献}

\begin{enumerate}
\item \textbf{元論文の再現と問題の特定}：公式実装のコードを読んで再現し、BLM がほぼビューだけになる理由（$\Omega$ が小さすぎる）、最適化の打ち切り、なぞりを数字で示した（第 VII・VIII 部）。
\item \textbf{評価の作り直し}：月次のリターンを当てる問いから、決算発表という\textbf{イベント}への反応を説明する問いに移した。株価の反応を計算する前に、計画とコードを固定する事前登録を毎回行った（第 X 部）。
\item \textbf{先読みの統制}：学習データの時点が決まっている ChronoGPT を使い、gpt-4o-mini の知識が及ばない期間でも、LLM の文章の読み取りが効くことを示した。
\item \textbf{比較の階層}：辞書・FinBERT・TF-IDF・追加学習した FinBERT・ChronoGPT の埋め込み・LoRA・gpt-4o-mini を同じ標本・同じ統制で比べた（\chapref{ch:conclusion}の表）。
\item \textbf{運用の道具の確認}：TE の約束が守られない問題を見つけ、過去の比で上限を校正する簡単な方法で直した（X03）。
\end{enumerate}
